% Paper 1 (theorems): T1, T3 (lead), T4 (pipeline validation), C1.
% Compile from paper/: pdflatex main_theorems.tex && pdflatex main_theorems.tex
% TITLE: P1-B (author-approved 2026-08-16)

\documentclass[onecolumn]{IEEEtran}
% Shared preamble fragments for the Karlsson MUB split papers.
% Usage: after \documentclass{...}, % Shared preamble fragments for the Karlsson MUB split papers.
% Usage: after \documentclass{...}, % Shared preamble fragments for the Karlsson MUB split papers.
% Usage: after \documentclass{...}, \input{preamble_common}
% Does not open \begin{document}.

\usepackage[T1]{fontenc}
\usepackage[utf8]{inputenc}
\usepackage{amsmath,amssymb,amsfonts}
\usepackage{graphicx}
\usepackage{cite}
\usepackage{booktabs}
\usepackage{array}
\usepackage{url}
\usepackage{xurl}
\usepackage{amsthm}
\usepackage[strings]{underscore}
\usepackage{hyperref}

% Allow line breaks in long repository paths and script names
\newcommand{\repo}[1]{{\ttfamily\nolinkurl{#1}}}
\emergencystretch=2.5em
\sloppy

\newtheorem{theorem}{Theorem}
\newtheorem{lemma}[theorem]{Lemma}
\newtheorem{finding}{Finding}
\newtheorem{conjecture}{Conjecture}
\newtheorem{remark}{Remark}

\hypersetup{
    colorlinks=true,
    linkcolor=black,
    citecolor=black,
    urlcolor=blue
}

\def\BibTeX{{\rm B\kern-.05em{\sc i\kern-.025em b}\kern-.08em
    T\kern-.1667em\lower.7ex\hbox{E}\kern-.125emX}}

% Does not open \begin{document}.

\usepackage[T1]{fontenc}
\usepackage[utf8]{inputenc}
\usepackage{amsmath,amssymb,amsfonts}
\usepackage{graphicx}
\usepackage{cite}
\usepackage{booktabs}
\usepackage{array}
\usepackage{url}
\usepackage{xurl}
\usepackage{amsthm}
\usepackage[strings]{underscore}
\usepackage{hyperref}

% Allow line breaks in long repository paths and script names
\newcommand{\repo}[1]{{\ttfamily\nolinkurl{#1}}}
\emergencystretch=2.5em
\sloppy

\newtheorem{theorem}{Theorem}
\newtheorem{lemma}[theorem]{Lemma}
\newtheorem{finding}{Finding}
\newtheorem{conjecture}{Conjecture}
\newtheorem{remark}{Remark}

\hypersetup{
    colorlinks=true,
    linkcolor=black,
    citecolor=black,
    urlcolor=blue
}

\def\BibTeX{{\rm B\kern-.05em{\sc i\kern-.025em b}\kern-.08em
    T\kern-.1667em\lower.7ex\hbox{E}\kern-.125emX}}

% Does not open \begin{document}.

\usepackage[T1]{fontenc}
\usepackage[utf8]{inputenc}
\usepackage{amsmath,amssymb,amsfonts}
\usepackage{graphicx}
\usepackage{cite}
\usepackage{booktabs}
\usepackage{array}
\usepackage{url}
\usepackage{xurl}
\usepackage{amsthm}
\usepackage[strings]{underscore}
\usepackage{hyperref}

% Allow line breaks in long repository paths and script names
\newcommand{\repo}[1]{{\ttfamily\nolinkurl{#1}}}
\emergencystretch=2.5em
\sloppy

\newtheorem{theorem}{Theorem}
\newtheorem{lemma}[theorem]{Lemma}
\newtheorem{finding}{Finding}
\newtheorem{conjecture}{Conjecture}
\newtheorem{remark}{Remark}

\hypersetup{
    colorlinks=true,
    linkcolor=black,
    citecolor=black,
    urlcolor=blue
}

\def\BibTeX{{\rm B\kern-.05em{\sc i\kern-.025em b}\kern-.08em
    T\kern-.1667em\lower.7ex\hbox{E}\kern-.125emX}}


% Convention (this document): Theorem/Lemma/Proof for T1 (algebra), T3
% (certified W_1 emptiness), and T4 (exact Groebner pipeline validation).
% Computational locus audits live in the companion methods report~\cite{JaupiMethods2026}.

\begin{document}

\title{Gauge Structure and Certified Non-Extendability for Selected Karlsson Hadamards in Dimension Six}

\author{%
\IEEEauthorblockN{Adian Jaupi}%
}

\maketitle

\begin{abstract}
Within Karlsson's three-parameter family $K_6^{(3)}$ only, I prove a gauge-structure theorem, give a certified fourth-MUB obstruction at seven Dita-$\lambda$ points, independently re-verify one Brierley--Weigert instance by exact Gr\"obner methods, and correct a literature transcription of Karlsson's $A$-matrix.
\textbf{(T1, Theorem~\ref{thm:T1}):} at $\theta{=}0$, $\phi$ is a parametrization gauge; on the Dita slice, $\lambda$ is not a matrix-level gauge, while CHM class is $\pi$-periodic (four classes among $\{0,0.4,\pi/3,2\pi/3,\pi,4\pi/3,5\pi/3\}$).
\textbf{(T3, Theorem~\ref{thm:T3}; lead result):} at seven Dita-$\lambda$ points $\lambda\in\{0,0.4,\pi/3,2\pi/3,\pi,4\pi/3,5\pi/3\}$, HomotopyContinuation \texttt{certify()} yields a numerical-interval certificate that the single-vector witness $W_1$ is empty at every recovered-pool $6$-clique ($\mathbf{40/40}$), which implies no fourth MUB at those instances.
\textbf{(T4, Theorem~\ref{thm:T4}; pipeline validation):} at the $D_0$-equivalent pair $(D_{\mathrm{bc}},F_D)$, an exact ideal-membership proof shows that the witness ideal $W_1$ is the unit ideal over $\mathbb{Q}(\zeta_{24},\sqrt5)$ (Gr\"obner basis $\{1\}$). This independently re-checks a fact already in Brierley--Weigert~\cite{Brierley2009}; it is not billed as a co-equal headline with T1/T3. A regression rerun of the pipeline (2026-09) additionally reproduces Brierley--Weigert's complete-pool counts at the CHM-equivalent point $\lambda{=}\pi/2$---$120$ vectors, $10$ third bases, no mutually unbiased pair among them, and all ten $W_1$ witnesses certified empty---so the exact certificate and the numerical pipeline now agree on the same instance. The closed-form third basis $F_D$ follows Bengtsson--Bruzda--Ericsson--Larsson--Tadej--\.Zyczkowski~\cite{BengtssonMUH2007}.
\textbf{(C1)} The McNulty--Weigert printed $A$-matrix differs from Karlsson's by a leading-sign error and a broken conjugate structure; symbolically $A_{\mathrm{MW}}A_{\mathrm{MW}}^\dagger\neq2I$ on a Zariski-open set (corroborated by a $0/349$ numerical audit).
Third-MUB locus identification, search-set audits, and reproducibility artifacts are documented in a companion computational report~\cite{JaupiMethods2026}.
\end{abstract}

\begin{IEEEkeywords}
Complex Hadamard matrices, Karlsson family, mutually unbiased bases, homotopy continuation, Gr\"obner bases, certified computation
\end{IEEEkeywords}

\section{Introduction}
\label{sec:intro}

Mutually unbiased bases (MUBs) in $\mathbb{C}^d$ are orthonormal bases $\mathcal{B}_0,\ldots,\mathcal{B}_{m-1}$ such that for all $a\neq b$ and all $|\psi\rangle\in\mathcal{B}_a$, $|\phi\rangle\in\mathcal{B}_b$,
\begin{equation}
    \label{eq:mu}
    \bigl|\langle\psi|\phi\rangle\bigr|^2 = \frac{1}{d}.
\end{equation}
MUBs play a central role in quantum state tomography, cryptography, and the study of finite geometry~\cite{Bengtsson2007, Brierley2009}.
For prime-power dimensions $d=p^n$, exactly $d+1$ MUBs exist~\cite{Ivanovic1981, Wootters1989}.
For composite $d$, the maximum number $N(d)$ of MUBs is unknown in general; in dimension six, Zauner conjectured $N(6)=3$~\cite{Zauner1999, Scott2010}.

A standard construction route passes through complex Hadamard matrices (CHMs): unimodular $d\times d$ matrices $H$ with $HH^\dagger=dI$.
Karlsson~\cite{Karlsson2011} parametrized a three-dimensional family $K_6(\theta,\phi,\lambda)$ of $6\times6$ CHMs.
Grassl~\cite{Grassl2004} showed that $\{\mathbf{I},F_6\}$ extends to at most three MUBs, so searches must be \emph{per} Hadamard matrix.
Brierley--Weigert~\cite{Brierley2009} constructed all MU vectors at the Di\c{t}a matrix $D_0$ in closed form and showed that none of the ten third bases are pairwise MU.

This paper is theorem-first within $K_6^{(3)}$ only.
The lead result is a certified fourth-MUB obstruction at seven Dita-$\lambda$ points (T3).
Supporting algebra is the gauge structure of Karlsson parameters (T1).
An exact Gr\"obner re-verification at one $D_0$-equivalent pair (T4) validates the witness pipeline against Brierley--Weigert but does not extend that literature.
A literature correction (C1) records that Karlsson's original $A$-matrix, not the McNulty--Weigert transcription, produces valid CHMs in numerical tests.
The computational identification of third-MUB loci (Dita circle; $F_6$ $\theta{=}0$ arc), the covering region $\mathcal{R}$, and reproducibility audits appear in the companion methods report~\cite{JaupiMethods2026}; the loci used to instantiate $B_3$ in the T3 runs are documented there.

\subsection{Contributions}
\begin{enumerate}
    \item[\textbf{T3}] \textbf{Lead:} certified $W_1$-emptiness (HomotopyContinuation \texttt{certify()}, tolerance-based) at seven Dita-$\lambda$ points, $40/40$ recovered-pool $6$-cliques (Theorem~\ref{thm:T3}).
    \item[\textbf{T1}] Gauge structure: $\phi$-gauge at $\theta{=}0$; $\lambda$-periodicity; Dita block specialization; Dita $\lambda$ matrix moduli with $\pi$-periodic CHM class (Theorem~\ref{thm:T1}).
    \item[\textbf{C1}] Literature correction: explicit algebraic discrepancy between Karlsson's $A$ and the McNulty--Weigert transcription (Section~\ref{sec:audit}).
    \item[\textbf{T4}] Pipeline validation: exact unit-ideal certificate for $W_1$ at $(D_{\mathrm{bc}},F_D)$, subordinate to Brierley--Weigert~\cite{Brierley2009} (Theorem~\ref{thm:T4}).
\end{enumerate}

\section{Scope and Logical Status of Results}
\label{sec:scope}

All results concern Karlsson's family $K_6(\theta,\phi,\lambda)$ only; CHMs outside $K_6^{(3)}$ are not treated.
\begin{itemize}
    \item \textbf{Proved (algebra).} Theorem~\ref{thm:T1} (T1); symbolic non-unitarity of the printed McNulty--Weigert $A$-block (C1).
    \item \textbf{Certified (numerical interval).} Theorem~\ref{thm:T3} (T3): \texttt{certify()} over a square subsystem of $W_1$ at seven listed points; not a certificate on the full $\lambda$-circle, on $\mathcal{R}$, or on $N(6)$.
    \item \textbf{Exact ideal membership (one point).} Theorem~\ref{thm:T4} (T4): Gr\"obner basis $\{1\}$ over $\mathbb{Q}(\zeta_{24},\sqrt5)$ at $(D_{\mathrm{bc}},F_D)$; independent re-verification of a BW fact, not a new extendability theorem on the circle.
    \item \textbf{Companion only (reproducibility).} Third-MUB locus identification, $\mathcal{S}^*$ sweeps, grid-miss narrative, pool-completeness tables, and the Open Problem on all of $\mathcal{R}$~\cite{JaupiMethods2026}. Theorem~\ref{thm:T3} does not logically depend on that Open Problem.
\end{itemize}
This paper does not prove Zauner's conjecture $N(6){=}3$ or fourth-MUB absence on all of $K_6^{(3)}$.

\section*{Main Theorem}
\label{sec:main-theorem}
\noindent\textbf{Lead result (Theorem~\ref{thm:T3}).}
For each
\[
\lambda\in\Lambda_{\mathrm{cert}}:=\Bigl\{0,\ 0.4,\ \tfrac{\pi}{3},\ \tfrac{2\pi}{3},\ \pi,\ \tfrac{4\pi}{3},\ \tfrac{5\pi}{3}\Bigr\}
\]
and every size-$6$ clique $B_3$ in the recovered MU pool of $\{\mathbf{I},H(\theta_D,\pi/4,\lambda)\}$ with $\theta_D=\arccos(1/\sqrt{3})$, the single-vector witness $W_1(H,B_3)$ is empty under HomotopyContinuation interval \texttt{certify()} of a generic square subsystem (mixed volume $252$), so $\{\mathbf{I},H,B_3\}$ admits no fourth MUB.

\noindent
\fbox{\parbox{0.92\linewidth}{\centering\vspace{4pt}
$\displaystyle\mathbf{40/40}$ certified obstructions
(native all-pool-clique exhaustion at the seven listed $\lambda$).\vspace{4pt}}}

\noindent Certificate details---including the square-subsystem completeness bridge ($N_{\mathrm{expected}}=N_{\mathrm{tracked}}=N_{\mathrm{certified}}=252$, $N_{\mathrm{full\text{-}system}}=0$)---are in Appendix~\ref{sec:proof-t3}; they are not deferred to the companion report.

\section{Background}
\label{sec:background}

\subsection{MUBs From Hadamard Matrices}
Let $H$ be a $6\times6$ CHM with $|H_{ij}|=1$ and $HH^\dagger=6I$.
The normalized columns $c_k=H_{:,k}/\sqrt{6}$ form an orthonormal basis $\mathcal{B}_H$ satisfying~\eqref{eq:mu} with the computational basis $\mathcal{B}_I=\{e_0,\ldots,e_5\}$.
A third MUB requires six unit vectors pairwise orthogonal and unbiased to both $\mathcal{B}_I$ and $\mathcal{B}_H$; a fourth MUB requires a second disjoint six-vector set with the same property relative to that third basis.

\subsection{Karlsson's Three-Parameter Family}
Following Karlsson~\cite{Karlsson2011} and McNulty--Weigert~\cite{McNulty2026},
\begin{equation}
    K_6(\theta,\phi,\lambda)=
    \begin{bmatrix}
        F_2 & Z_1 & Z_2 \\
        Z_3 & \tfrac{1}{2}Z_3 A Z_1 & \tfrac{1}{2}Z_3 B Z_2 \\
        Z_4 & \tfrac{1}{2}Z_4 B Z_1 & \tfrac{1}{2}Z_4 A Z_2
    \end{bmatrix},
\end{equation}
where $F_2=\begin{pmatrix}1&1\\1&-1\end{pmatrix}$, $B=-F_2-A$, and $Z_j$ are $2\times2$ blocks from unimodular parameters $z_1,\ldots,z_4$ with $z_1=e^{i\lambda}$.
The block $A$ depends on $(\theta,\phi)$.
Karlsson's original $A$-matrix formulas are used throughout; the McNulty--Weigert printed transcription is algebraically invalid as a CHM block (Section~\ref{sec:audit}; contribution C1).

\section{Methods}
\label{sec:methods}

\subsection{Per-$H$ MU pool and fourth-MUB witness}
For fixed $H$, seek unit vectors unbiased to $\mathcal{B}_I$ and $\mathcal{B}_H$.
With gauge $v_0=1$ and $z_i=v_i\sqrt{6}$, unbiasedness to the five independent columns of $H$ together with unimodular auxiliaries $z_iw_i=1$ yields a square $10\times10$ polynomial system depending on $H$.
Solutions on the conjugate locus are mapped to unit vectors, verified at tolerance $10^{-8}$, and deduplicated modulo global phase by projective (Fubini--Study) clustering at $128$-bit precision; an independent re-seeded re-solve must agree on the physical count before a pool is accepted, and the resulting finite set is the recovered MU pool $\mathcal{P}$.
A third basis $B_3$ is any orthonormal $6$-frame whose columns form a size-$6$ clique in the orthogonality graph on $\mathcal{P}$ (edge when $|\langle v,w\rangle|<\varepsilon_{\mathrm{ortho}}$).
The single-vector witness $W_1(H,B_3)$ encodes existence of a unit vector unbiased to $\mathbf{I}$, $H$, and $B_3$ ($17$ equations in $10$ variables after unimodular auxiliaries; $n_{\mathrm{wit}}=1$).
Emptiness of $W_1$ implies no fourth MUB at that $(H,B_3)$.
(Full pipeline engineering, search sets, and reproducibility pins appear in~\cite{JaupiMethods2026}; they are not required to read the T3 certificate.)

\subsection{T3 certificate (numerical interval)}
At each listed $\lambda$, every recovered-pool $6$-clique is submitted to a generic square subsystem of $W_1$ (mixed volume $252$), HomotopyContinuation path tracking, interval-arithmetic \texttt{certify()}, and residual check against the full system (tolerance $10^{-8}$).
Emptiness means $N_{\mathrm{expected}}=N_{\mathrm{tracked}}=N_{\mathrm{certified}}=252$ and $N_{\mathrm{full\text{-}system}}=0$ (Appendix~\ref{sec:proof-t3}).
This is \emph{not} a Gr\"obner unit-ideal certificate.

\subsection{T4 certificate (exact ideal membership)}
At $(D_{\mathrm{bc}},F_D)$ with algebraic entries in $\mathbb{Q}(\zeta_{24},\sqrt5)$, Macaulay2 computes a Gr\"obner basis of the $W_1$ ideal equal to $\{1\}$ (Section~\ref{sec:proof-t3}).

\subsection{Central open obstacle}
\label{sec:symbolic}
A theorem of fourth-MUB absence on all of a covering region $\mathcal{R}\subset K_6^{(3)}$ (defined in~\cite{JaupiMethods2026}) would require exact-coefficient elimination with $B_3$ in a genuine number field at points not CHM-equivalent to $D_0$.
The reduced two-vector witness run over \texttt{CC} is inconclusive (inexact-field warning; Remark~\ref{rem:m2-inexact}).
T4 works precisely because $B_3=F_D$ is closed-form at a $D_0$-equivalent point---corollary-grade relative to Brierley--Weigert as a mathematical result, valuable here as a machine-checkable cross-check.
Theorem~\ref{thm:T3} does not depend on resolving that open obstacle.

\section{Results}
\label{sec:results}

\subsection{Correction to the McNulty--Weigert \texorpdfstring{$A$}{A}-matrix transcription}
\label{sec:audit}

Karlsson~\cite{Karlsson2011} uses
\[
A_{11}=-\tfrac12+i\tfrac{\sqrt3}{2}\bigl(\cos\theta+e^{-i\phi}\sin\theta\bigr),\quad
A_{12}=-\tfrac12+i\tfrac{\sqrt3}{2}\bigl(-\cos\theta+e^{i\phi}\sin\theta\bigr),
\]
with Hermitian block structure $A=\begin{pmatrix}A_{11}&A_{12}\\ \overline{A_{12}}&-\overline{A_{11}}\end{pmatrix}$.
Eqs.~(7.3)--(7.5) of~\cite{McNulty2026}, as printed, instead take leading $+\tfrac12$ in $A_{11}$ and assemble $A=\begin{pmatrix}A_{11}&A_{12}\\ A_{12}&-A_{11}\end{pmatrix}$ (no conjugate on the off-diagonal).

\begin{lemma}[Printed McNulty--Weigert $A$ is not unitary]
\label{lem:C1-symbolic}
Let $A_{\mathrm{MW}}(\theta,\phi)$ be the printed block above.
Then
\[
\bigl(A_{\mathrm{MW}}A_{\mathrm{MW}}^\dagger-2I\bigr)_{11}
=\sqrt3\,\sin\phi\,\sin\theta,
\]
which is not identically zero.
Hence $A_{\mathrm{MW}}A_{\mathrm{MW}}^\dagger\neq2I$ on a Zariski-open subset of $(\theta,\phi)$, so the printed transcription cannot define a valid Karlsson CHM on that open set.
Moreover $A_{11}^{\mathrm{(Karlsson)}}-A_{11}^{\mathrm{(MW)}}=-1$ identically (leading-sign discrepancy).
\end{lemma}

\begin{proof}
Direct expansion of $A_{\mathrm{MW}}A_{\mathrm{MW}}^\dagger$ (computer algebra; reproducible from \repo{scripts/python/audit\_karlsson\_variants.py} formulas) yields the displayed $(1,1)$ entry.
The leading-sign identity is immediate from $\pm\tfrac12$.
\end{proof}

\begin{table}[t]
\caption{Correction summary for the $A$-matrix transcription.}
\label{tab:audit-summary}
\centering
\begin{tabular}{@{}lcc@{}}
\toprule
 & Karlsson 2011 & McNulty--Weigert (as printed) \\
\midrule
Symbolic $A$ formula & valid & leading $+1/2$; non-conjugate structure \\
$AA^\dagger=2I$ & identity & fails: $(AA^\dagger-2I)_{11}=\sqrt3\sin\phi\sin\theta$ \\
$349$-point numerical audit & $348/349$ pass & $0/349$ pass \\
Identified error & --- & Eqs.~(7.3)--(7.5) as printed \\
\bottomrule
\end{tabular}
\end{table}

\begin{table}[t]
\caption{Numerical audit of $A$-matrix variants on $349$ parameter points (corroboration of Lemma~\ref{lem:C1-symbolic}). Pass criteria: $\|AA^\dagger-2I\|_\infty<10^{-8}$, $\|BB^\dagger-2I\|_\infty<10^{-8}$, valid CHM, M\"obius consistency $<10^{-6}$.}
\label{tab:audit}
\centering
\begin{tabular}{@{}lccc@{}}
\toprule
Variant & Pass & Fail & Notes \\
\midrule
Karlsson 2011 (original) & 348 & 1 & One NaN at degenerate M\"obius point \\
McNulty--Weigert (as printed) & 0 & 349 & $A$ not unitary; CHM invalid \\
\bottomrule
\end{tabular}
\end{table}

\noindent
The correction is also enforced as a regression test: the literal transcription fails $AA^\dagger=2I$ at $40/40$ random parameter points in the repository test suite, while Karlsson's original passes at every point where the family is defined.

\subsection{Gauge Structure (T1)}
Theorem~\ref{thm:T1} records that at $\theta{=}0$, $\phi$ does not enter $H$; the family is $2\pi$-periodic in $\lambda$; on the Dita slice $\lambda$ changes the matrix (Lemma~\ref{lem:L4-matrix}) while CHM class is $\pi$-periodic with four classes among the seven sample points of Lemma~\ref{lem:L4}.
Proof: Appendix~\ref{sec:proofs}.

\subsection{Certified Fourth-MUB Obstruction at Seven Dita Points (T3)}
\label{sec:t3-results}
Theorem~\ref{thm:T3} is the lead result (Section~\ref{sec:main-theorem}): at
$\Lambda_{\mathrm{cert}}=\{0,0.4,\pi/3,2\pi/3,\pi,4\pi/3,5\pi/3\}$,
\texttt{certify()} finds $W_1$ empty for all recovered-pool $6$-cliques ($\mathbf{40/40}$), covering four CHM classes (Lemma~\ref{lem:L4}).
Full proof: Appendix~\ref{sec:proofs}.
Confirmatory dense probes and search-set negatives outside these seven points are reproducibility material in~\cite{JaupiMethods2026}; they are not part of T3.

\subsection{Pipeline Validation: Exact Re-Verification at the \texorpdfstring{$D_0$}{D0}-Equivalent Pair (T4)}
\label{sec:t4-results}
Theorem~\ref{thm:T4} establishes $W_1(D_{\mathrm{bc}},F_D)=\langle 1\rangle$ by Gr\"obner basis over $\mathbb{Q}(\zeta_{24},\sqrt5)$.
Brierley--Weigert already proved the stronger complete-pool statement ($N_v{=}120$, $N_t{=}10$, $N_p{=}0$); T4 covers one third basis by a different exact method and is billed here as independent pipeline validation (Remark~\ref{rem:T4-scope}).
After the methodology hardening documented in the companion report~\cite{JaupiMethods2026}, the same pipeline also \emph{reproduces} the complete-pool statement numerically at the CHM-equivalent point $\lambda{=}\pi/2$ on the Dita slice: a certified solve returns exactly $120$ physical MU vectors ($N_v$) forming $10$ independently verified third bases ($N_t$), with no mutually unbiased pair among all $45$ basis pairs ($N_p{=}0$; worst defect $0.015$), and the single-vector witness certified empty for all $10/10$ cliques---the four Brierley--Weigert counts, regenerated end-to-end by the audited pipeline rather than imported from the literature.

\section{Discussion}
\label{sec:discussion}

T3 is a pointwise certified obstruction, not a theorem on $\mathcal{R}$ or $N(6)$.
T4 does not extend Brierley--Weigert; it cross-checks one algebraic instance of the same witness ideal used numerically in T3.
The September 2026 regression rerun (pool counts at $F_6$, the $D_0$-equivalent $\lambda{=}\pi/2$ point, and Tao's isolated order-$6$ matrix) is a pipeline-audit artifact documented with the companion report~\cite{JaupiMethods2026}; it changes no claim tier in this paper.
C1 is an algebraic transcription correction (Lemma~\ref{lem:C1-symbolic}) and should be reported to the authors of~\cite{McNulty2026}.
The companion report~\cite{JaupiMethods2026} documents search-set reproducibility and why dense-grid negatives alone cannot promote sweep results to theorems---reinforcing that T3 is the exception that does not rest on grid density.

\section{Conclusion}
\label{sec:conclusion}

Scoped to $K_6^{(3)}$ only, this paper delivers a gauge theorem (T1), a lead certified fourth-MUB obstruction at seven Dita-$\lambda$ points (T3), an algebraic literature correction (C1), and an exact pipeline-validation re-check at one $D_0$-equivalent pair (T4), subordinate to Brierley--Weigert.
A post-audit regression run additionally reproduces Brierley--Weigert's complete-pool counts end-to-end ($120$ vectors, $10$ third bases, $N_p{=}0$, $W_1$ empty $10/10$) at the CHM-equivalent point, closing the loop between the exact certificate and the numerical pipeline.
I do not claim fourth-MUB absence on all of $K_6^{(3)}$ or on $N(6)$.
Closing the gap to a theorem on $\mathcal{R}$ requires algebraic $B_3$ at points not CHM-equivalent to $D_0$; see~\cite{JaupiMethods2026} for the reproducibility record of that open obstacle.

\section*{Acknowledgment}
I thank the open-source developers of HomotopyContinuation.jl and the maintainers of the Karlsson CHM literature surveys.

\appendices
\section{Proofs of T1, T3, and T4}
\label{sec:proofs}
Theorem and Proof are used for T1 (Lemmas~\ref{lem:L1}--\ref{lem:L4}), T3, and T4.
% !TeX root = ../main_theorems_multifile.tex
% Paper 1 appendix: algebraic gauge structure (T1). Locus geometry lives in Paper 2.

\section{Gauge structure in Karlsson's family}
\label{sec:proof-gauge}

Write Karlsson's three-parameter family as $K_6(\theta,\phi,\lambda)$ with block $A(\theta,\phi)$ and $z_1=e^{i\lambda}$ determining the remaining M\"obius blocks (as in Section~\ref{sec:background}).
Throughout, $H(\theta,\phi,\lambda):=K_6(\theta,\phi,\lambda)$.

\begin{lemma}[Block $A$ at $\theta=0$]
\label{lem:phi-gauge-A}
For all $\phi,\phi'\in[0,\pi)$ and $\lambda\in\mathbb{R}$,
\[
A(0,\phi)=A(0,\phi')=
\begin{bmatrix}
-\tfrac12 + i\tfrac{\sqrt3}{2} & -\tfrac12 - i\tfrac{\sqrt3}{2} \\
-\tfrac12 + i\tfrac{\sqrt3}{2} & \tfrac12 + i\tfrac{\sqrt3}{2}
\end{bmatrix}.
\]
\end{lemma}

\begin{proof}
Substitute $\sin\theta=0$ in the Karlsson formulas
\[
A_{11}=-\tfrac12 + i\tfrac{\sqrt3}{2}(\cos\theta+e^{-i\phi}\sin\theta),\quad
A_{12}=-\tfrac12 + i\tfrac{\sqrt3}{2}(-\cos\theta+e^{i\phi}\sin\theta).
\]
All $\phi$-dependent terms vanish, giving the displayed matrix.
Unitarity $AA^\dagger=2I$ follows by direct multiplication.
\end{proof}

\begin{lemma}[$\phi$-gauge for $H$ at $\theta=0$]
\label{lem:L1}
For all $\phi,\phi',\lambda$: $H(0,\phi,\lambda)=H(0,\phi',\lambda)$ entrywise.
\end{lemma}

\begin{proof}
At $\theta=0$, Lemma~\ref{lem:phi-gauge-A} gives $A$ and $B=-F_2-A$ independent of $\phi$.
The M\"obius parameters $\alpha_A,\beta_A,\alpha_B,\beta_B$ depend only on $A,B$, hence on $\theta$ alone.
With $z_1=e^{i\lambda}$, the quantities $z_2,z_3,z_4$ and all blocks $Z_j$ are $\phi$-independent.
Therefore the assembled $6\times6$ matrix $H(0,\phi,\lambda)$ is independent of $\phi$.
\end{proof}

\begin{lemma}[$\lambda$-periodicity]
\label{lem:L2}
For all $\theta,\phi,\lambda$: $H(\theta,\phi,\lambda+2\pi)=H(\theta,\phi,\lambda)$.
\end{lemma}

\begin{proof}
The only $\lambda$-dependence enters through $z_1=e^{i\lambda}$, which satisfies $e^{i(\lambda+2\pi)}=e^{i\lambda}$.
All subsequent M\"obius/square-root steps depend on $z_1^2$ and hence are $2\pi$-periodic in $\lambda$.
\end{proof}

\begin{lemma}[Dita block specialization]
\label{lem:L3}
At $\theta=\arccos(1/\sqrt3)$ and $\phi=\pi/4$,
\[
A(\theta_D,\pi/4)=
\begin{bmatrix} i & -1 \\ -1 & i \end{bmatrix},
\qquad AA^\dagger=2I.
\]
\end{lemma}

\begin{proof}
With $\cos\theta=1/\sqrt3$, $\sin\theta=\sqrt{2/3}$, and $\phi=\pi/4$,
\[
A_{11}=-\tfrac12+i\tfrac{\sqrt3}{2}\Bigl(\tfrac{1}{\sqrt3}+\tfrac{1-i}{\sqrt3}\Bigr)=i,\quad
A_{12}=-\tfrac12+i\tfrac{\sqrt3}{2}\Bigl(-\tfrac{1}{\sqrt3}+\tfrac{1+i}{\sqrt3}\Bigr)=-1.
\]
Hermitian structure $A=\begin{pmatrix}A_{11}&A_{12}\\ \overline{A_{12}}&-\overline{A_{11}}\end{pmatrix}$ yields the matrix shown.
Direct computation gives $AA^\dagger=2I$.
\end{proof}

\begin{lemma}[Matrix inequality in $\lambda$ at Dita]
\label{lem:L4-matrix}
Let $\theta_D=\arccos(1/\sqrt3)$, $\phi=\pi/4$, and write $H(\lambda):=H(\theta_D,\phi,\lambda)$.
If $\lambda\not\equiv\lambda'\pmod{2\pi}$, then $H(\lambda)\neq H(\lambda')$ as matrices.
\end{lemma}

\begin{proof}
In the block assembly, the one-based entry $H(\lambda)_{2,3}$ is $(Z_1)_{2,1}=z_1=e^{i\lambda}$. Therefore equality of $H(\lambda)$ and $H(\lambda')$ implies $e^{i\lambda}=e^{i\lambda'}$, hence $\lambda\equiv\lambda'\pmod{2\pi}$.
\end{proof}

\begin{lemma}[CHM class structure on the Dita $\lambda$-circle]
\label{lem:L4}
\label{lem:dita-chm-pi}
Let $\theta_D=\arccos(1/\sqrt3)$, $\phi=\pi/4$, and $H(\lambda):=H(\theta_D,\phi,\lambda)$.
\begin{enumerate}
    \item \emph{$\pi$-periodicity of CHM class.} Dephased alignment over $720$ column permutations and diagonal phases yields residuals $\sim10^{-15}$ for the pairs $(\lambda,\lambda+\pi)$ with $\lambda\in\{0,\pi/3,2\pi/3\}$ (\repo{chm\_equivalence\_residual}).
    Thus $H(\lambda)$ and $H(\lambda+\pi)$ are CHM-equivalent on this slice.
    \item \emph{Four distinct classes at standard sample points.} The representatives $\{H(0),H(0.4),H(\pi/3),H(2\pi/3)\}$ are pairwise CHM-inequivalent (minimum residual $\gtrsim0.49$; \repo{results/verify\_nwit1\_referee\_checks.txt}).
    Combined with item~(1), the seven values $\{0,0.4,\pi/3,2\pi/3,\pi,4\pi/3,5\pi/3\}$ cover exactly these four CHM classes.
    \item \emph{Not a parametrization gauge.} Distinct $\lambda$ modulo $2\pi$ never yield identical matrices (Lemma~\ref{lem:L4-matrix}); $\lambda$ is a genuine moduli coordinate at the matrix level, even though CHM class is $\pi$-periodic.
\end{enumerate}
In particular, it is \textbf{not} claimed that every pair $\lambda\not\equiv\lambda'\pmod{2\pi}$ gives CHM-inequivalent matrices---that would contradict item~(1).
A complete invariant-theoretic classification of all $\lambda\in[0,2\pi)$ modulo diagonal-phase/permutation CHM moves remains open; items~(1)--(2) record the numerical class structure used for the seven-point endpoint tests below.
\end{lemma}

\begin{proof}
Item~(3) is Lemma~\ref{lem:L4-matrix}.
Items~(1)--(2) are the CHM-equivalence computations recorded in the cited verification logs.
\end{proof}

\begin{theorem}[Gauge structure, T1]
\label{thm:T1}
Within Karlsson's family $K_6^{(3)}$:
\begin{enumerate}
    \item At $\theta=0$, the coordinate $\phi$ is a parametrization gauge: $H(0,\phi,\lambda)$ is constant in $\phi$ (Lemma~\ref{lem:L1}).
    \item The family is $2\pi$-periodic in $\lambda$ (Lemma~\ref{lem:L2}).
    \item At the Dita slice $(\theta_D,\pi/4)$, the block $A$ is the fixed unitary in Lemma~\ref{lem:L3}.
    \item On the Dita slice, $\lambda$ is not a matrix-level gauge (Lemma~\ref{lem:L4-matrix}), while CHM class is $\pi$-periodic with four distinct classes among $\{0,0.4,\pi/3,2\pi/3,\pi,4\pi/3,5\pi/3\}$ (Lemma~\ref{lem:L4}).
\end{enumerate}
\end{theorem}

\begin{proof}
Immediate from Lemmas~\ref{lem:L1}--\ref{lem:L4}.
\end{proof}

% !TeX root = ../main_theorems_multifile.tex
% Paper 1 appendix: T3 (certified) and T4 (exact re-verification). Finding L7 moved to methods companion.

\section{Fourth-MUB obstruction on the Dita slice (T3)}
\label{sec:proof-t3}

\begin{lemma}[CHM equivalence transfers $W_1$-emptiness]\label{lem:chm-w1}
Suppose $H_1,H_2\in K_6^{(3)}$ are CHM-equivalent via diagonal unitaries $D_r,D_c$ and a column permutation $P$, in the sense that after the standard dephasing gauge
\[
\mathrm{dephase}(H_1)=D_r\,\mathrm{dephase}(H_2)\,P\,D_c.
\]
Let $B_3$ be an orthonormal $6\times6$ matrix whose columns are mutually unbiased to $\mathbf{I}$ and to the columns of $H_2$, and set $B_3':=D_r B_3$ (columnwise, then renormalized).
Then:
\begin{enumerate}
    \item the columns of $B_3'$ are orthonormal and mutually unbiased to $\mathbf{I}$ and to the columns of $H_1$ (hence $B_3'$ is a third MUB for $\{\mathbf{I},H_1\}$);
    \item a unit vector $v$ solves $W_1(H_2,B_3)$ if and only if $D_r v$ solves $W_1(H_1,B_3')$.
\end{enumerate}
In particular, $W_1(H_2,B_3)$ is empty if and only if $W_1(H_1,B_3')$ is empty.
\end{lemma}

\begin{proof}[Sketch]
Diagonal unitaries act entrywise by phases, so $| (D_r v)_i |=|v_i|$ and $\langle D_r u, D_r w\rangle=\langle u,w\rangle$.
Mutual unbiasedness to $\mathbf{I}$ and pairwise orthogonality are therefore preserved by $v\mapsto D_r v$ and $B_3\mapsto B_3'$.
For mutual unbiasedness to the Hadamard: the CHM relation identifies columns of (dephased) $H_1$ with phased, permuted columns of (dephased) $H_2$ left-multiplied by $D_r$, so
\[
\bigl|\langle D_r v,\, (H_1)_{:,k}\rangle\bigr|
=
\bigl|\langle v,\, (H_2)_{:,p(k)}\rangle\bigr|
\]
up to the dephasing gauge (itself diagonal phases).
Thus MU-to-$H$ and MU-to-$B_3$ pass to MU-to-$H_1$ and MU-to-$B_3'$, and conversely.
\end{proof}

\begin{theorem}[No fourth MUB at certified Dita-$\lambda$ points]
\label{thm:T3}
\label{claim:T3}
Let $\theta_D=\arccos(1/\sqrt{3})$, $\phi=\pi/4$, and
\[
\Lambda_{\mathrm{cert}}:=\Bigl\{0,\ 0.4,\ \tfrac{\pi}{3},\ \tfrac{2\pi}{3},\ \pi,\ \tfrac{4\pi}{3},\ \tfrac{5\pi}{3}\Bigr\}.
\]
Write $H(\lambda)=H(\theta_D,\phi,\lambda)$.
For every $\lambda\in\Lambda_{\mathrm{cert}}$ and every third MUB $B_3$ arising as a size-$6$ clique in the recovered MU pool of $\{\mathbf{I},H(\lambda)\}$ (\S\ref{sec:methods}), the triple $\{\mathbf{I},H(\lambda),B_3\}$ does \emph{not} extend to a fourth mutually unbiased basis.
\end{theorem}

\begin{proof}
\textbf{Logic.}
A fourth MUB $B_4=\{v_1,\ldots,v_6\}$ extending $\{\mathbf{I},H,B_3\}$ consists of six orthonormal vectors, each mutually unbiased to every column of $\mathbf{I}$, $H$, and $B_3$.
After a global phase fixing the first coordinate positive, each $v_j$ is a solution of the single-vector witness $W_1(H,B_3)$:
five unimodular gauge equations $|z_i|=1$ ($i=1,\ldots,5$) together with mutual-unbiasedness to the six columns of $H$ and of $B_3$ ($17$ equations in $10$ complex variables; \repo{build\_numeric\_fourth\_mub\_witness\_system}, $n_{\mathrm{wit}}=1$).
Hence a fourth MUB implies feasibility of $W_1$; contrapositively, emptiness of $W_1$ proves absence of a fourth MUB at that $(H,B_3)$.\footnote{The converse fails in general: $W_1$ may be feasible without six mutually orthogonal solutions. Emptiness is therefore a complete---indeed strictly stronger---certificate of fourth-MUB non-existence at the point.}

\textbf{Method (every certified instance).}
Form a generic square subsystem of $W_1$ by ten random complex linear combinations of the seventeen equations (Bertini / HomotopyContinuation pattern; fixed seed), compute its mixed volume ($252$ in every run below), track all $252$ start paths, call HomotopyContinuation's interval-arithmetic \texttt{certify()} on every nonsingular endpoint, and residual-check each certified square-root against the full overdetermined system (tolerance $10^{-8}$).
For each submitted $(H,B_3)$ the accounting is
\[
N_{\mathrm{expected}}=N_{\mathrm{tracked}}=N_{\mathrm{certified}}=252,\qquad
N_{\mathrm{full\text{-}system}}:=\texttt{full\_residual\_pass}=0.
\]

\textbf{Square-subsystem completeness bridge.}
The full witness $W_1$ is overdetermined ($17$ equations in $10$ complex variables after unimodular auxiliaries).
A generic choice of ten complex linear combinations yields a square system whose isolated roots, by the standard Bertini / generic-linear-combination principle used in numerical algebraic geometry, include every isolated solution of the full system that remains isolated under that generic combination.
Path tracking recovers all $N_{\mathrm{expected}}=252$ start paths of the square system (mixed volume $252$ throughout); interval-arithmetic \texttt{certify()} validates nonsingular endpoints ($N_{\mathrm{certified}}=252$); residual filtering against the full $17$-equation system returns $N_{\mathrm{full\text{-}system}}=0$.
Hence no approximate root of the full witness appears among the certified square-system endpoints.
This is a numerical-interval certificate of emptiness for that $(H,B_3)$, \emph{not} a Gr\"obner unit-ideal membership proof (contrast Theorem~\ref{thm:T4}).

\textbf{Native all-pool-clique exhaustion at seven $\lambda$.}
At each $\lambda\in\Lambda_{\mathrm{cert}}$, every distinct size-$6$ clique in the recovered MU pool was submitted to the method above:
\[
\begin{array}{c|cc}
\lambda & \#\{6\text{-cliques}\} & \text{empty}\\
\hline
0 & 10 & 10/10\\
0.4 & 4 & 4/4\\
\pi/3 & 4 & 4/4\\
2\pi/3 & 4 & 4/4\\
\pi & 10 & 10/10\\
4\pi/3 & 4 & 4/4\\
5\pi/3 & 4 & 4/4
\end{array}
\]
Total: $\mathbf{40/40}$ empty (\repo{results/certify\_nwit1\_all\_cliques\_four\_classes.txt}, \repo{results/certify\_nwit1\_all\_cliques\_lambdapi.txt}).
In each run \texttt{found\_fourth=false} as well.

\textbf{CHM-class structure (independence, not inheritance).}
The $720$-permutation / diagonal-phase search (\repo{chm\_equivalence\_residual}) shows $H(\lambda)\sim H(\lambda+\pi)$ on this slice (residuals $\sim10^{-15}$), while
$\{H(0),H(0.4),H(\pi/3),H(2\pi/3)\}$ are pairwise CHM-inequivalent (minimum residual $\gtrsim0.49$; \repo{results/verify\_nwit1\_referee\_checks.txt}), as recorded in Lemma~\ref{lem:L4}.
The seven points therefore cover four CHM classes; the native $40/40$ check is independent at each $\lambda$ (not deduced from representatives).
As corroboration, CHM transport of every clique at $\{0,\pi/3,2\pi/3\}$ to $\{\pi,4\pi/3,5\pi/3\}$ also yields $18/18$ empty (\repo{results/verify\_chm\_b3\_transfer.txt}; Lemma~\ref{lem:chm-w1}).

\textbf{Clique-index nondeterminism.}
Among equal-length $6$-cliques, Graphs.jl enumeration order need not repeat across re-solves, so a single \texttt{argmax(|C|)} index set is not a stable label.
The native $40/40$ exhaustion removes dependence on that choice by ranging over every size-$6$ clique in the recovered pool at each $\lambda$.
\end{proof}

\begin{remark}[Scope --- what is and is not claimed]
Theorem~\ref{thm:T3} is a \emph{certified} non-existence result at seven explicit Dita-$\lambda$ values (four CHM classes), with native all-pool-clique coverage at every listed $\lambda$ ($40/40$).
``Every size-$6$ clique'' means every such clique in the MU pool recovered by the per-$H$ homotopy at that run (pool sizes $120$ at $\lambda{=}0,\pi$ and $72$ at the other five listed $\lambda$); it does not by itself certify against hypothetical third MUBs missed by an incomplete root track.
The $72$/$120$ split is a $\lambda$-dependent pool-size invariant of the pair $\{\mathbf{I},H(\lambda)\}$, not a pipeline artifact: the 2026-09 regression benchmarks (companion report) show the pool carries $120$ vectors in four $\lambda$-intervals of width ${\approx}0.22$ rad centered at $\lambda\in\{0,\pi/2,\pi,3\pi/2\}$ (mod $2\pi$) and $72$ vectors elsewhere on the circle, with $10$ third bases in the $120$-regions and $4$ in the $72$-regions, and with deduplicated counts stable under a $10^{-6}$--$10^{-12}$ dedup-tolerance sweep; the benchmark third-basis counts ($10$ at $\lambda{=}\pi/2$, $4$ at $\lambda{=}0.4$) match the clique counts in the proof above.
At the remaining $120$-point $\lambda{=}\pi/2$ (CHM-equivalent to $D_0$), the same benchmark certifies $W_1$ empty for all $10/10$ recovered cliques---numerically consistent with the pattern certified here at $\lambda{=}0,\pi$.
It is \textbf{not} a symbolic identity for all $\lambda\in[0,2\pi)$, not a theorem on the covering region $\mathcal{R}$ defined in the companion methods report (Paper~2), and not a statement about $N(6)$ or CHMs outside $K_6^{(3)}$.
A separate dense probe finds no fourth MUB at $628/628$ samples on $[0,2\pi]$; that numerical evidence is documented in Paper~2 and is \textbf{not} part of the certificate in Theorem~\ref{thm:T3}.
\end{remark}

\begin{remark}[Pointer to the central open obstacle]\label{rem:m2-inexact}
The two-vector Macaulay2 run over \texttt{CC} ($\dim I=-1$ with an inexact-field warning) is \textbf{not} a certificate and is \textbf{not} confirmatory of T3.
T3 uses $n_{\mathrm{wit}}=1$ (mixed volume $252$) instead.
The failed exact-elimination attempt is the central obstacle to a theorem on all of $\mathcal{R}$; see \S\ref{sec:symbolic} and the companion methods audit (Paper~2).
\end{remark}

\begin{theorem}[Exact-coefficient $W_1$ emptiness at the $D_0$-equivalent pair (T4)]
\label{thm:T4}
Let $D_{\mathrm{bc}}$ be the block-circulant representative of the Di\c{t}a matrix $D_0$ and $F_D$ the closed-form third MUB for $\{\mathbf{I},D_{\mathrm{bc}}\}$ (Bengtsson--Bruzda--Ericsson--Larsson--Tadej--\.Zyczkowski 2007, eqs.~(79)--(80) and (78),(61)), both with entries in $\mathbb{Q}(\zeta_{24},\sqrt{5})$.
Then the single-vector witness ideal of $W_1(D_{\mathrm{bc}},F_D)$ --- seventeen generators (five unimodular relations $z_iw_i=1$ and twelve mutual-unbiasedness products) in $\mathbb{Q}(\zeta_{24},\sqrt5)[z_1,\ldots,z_5,w_1,\ldots,w_5]$; \repo{symbolic\_export/w1\_D0\_exact.m2} --- is the \emph{unit ideal}.
Consequently no flat vector is mutually unbiased to $\mathbf{I}$, $D_{\mathrm{bc}}$, and $F_D$ simultaneously: the MUB triple $\{\mathbf{I},D_{\mathrm{bc}},F_D\}$ does not extend by even a single vector of a fourth MUB.
\end{theorem}

\begin{proof}
\textbf{Logical status.}
The implication is the same as in Theorem~\ref{thm:T3} ($W_1$ empty $\Rightarrow$ no fourth MUB at that $(H,B_3)$), but emptiness here is established as ideal-theoretic unit-ideal membership over an exact number field --- a Gr\"obner certificate, not a homotopy \texttt{certify()} run.

\textbf{Field validity.}
$[\mathbb{Q}(\zeta_{24},\sqrt5):\mathbb{Q}]=16$: $\Phi_{24}(x)=x^8-x^4+1$ is irreducible over $\mathbb{Q}$, and $5$ is not a square in $\mathbb{Q}(\zeta_{24})$ --- certified by the split prime $p=73$ ($p\equiv1\bmod 24$, so $\Phi_{24}$ splits completely mod $73$ and $\mathbb{Q}(\zeta_{24})$ embeds in residue fields $\mathbb{F}_{73}$, where $5$ is a quadratic non-residue).
Hence Macaulay2's \texttt{toField} quotient $\mathbb{Q}[\mathrm{zt},s_5]/(\Phi_{24}(\mathrm{zt}),\,s_5^2-5)$ is genuinely a field and Gr\"obner arithmetic over it involves no zero divisors (\repo{scripts/julia/verify\_field\_degree\_w1.jl}; \repo{results/track\_c\_elimination/w1\_D0\_field\_degree.txt}).

\textbf{Method and verification.}
Macaulay2~1.24.05, bounded cost probe first, then the certificate run:
\[
\begin{array}{l|l|l|r}
\text{run} & \text{field} & \text{result} & \text{wall}\\
\hline
\text{probe (}\zeta_{24}\mapsto2,\ \sqrt5\mapsto103\text{)} & \mathbb{Z}/241 & \text{unit ideal (signal only)} & 0.12\,\mathrm{s}\\
\text{bounded \texttt{dim}} & \mathbb{Q}(\zeta_{24},\sqrt5) & \dim = -1 & 29.0\,\mathrm{s}\\
\text{certificate} & \mathbb{Q}(\zeta_{24},\sqrt5) & \mathrm{GB}=\{1\},\ 1\,\%\,I = 0 & 16.0\,\mathrm{s}
\end{array}
\]
The certificate run computes the full Gr\"obner basis (\texttt{gens gb}), which is the single element $1$, and confirms $1\,\%\,I=0$ (\repo{symbolic\_export/w1\_D0\_groebner.m2}; \repo{results/track\_c\_elimination/w1\_D0\_groebner.log}, \repo{w1\_D0\_cost\_estimate.log}).
No inexact-field warning arises (contrast Remark~\ref{rem:m2-inexact}).
\end{proof}

\begin{remark}[Scope of Theorem~\ref{thm:T4} --- what is and is not claimed]
\label{rem:T4-scope}
(i) \emph{Point scope.} The exact statement is at the pair $(D_{\mathrm{bc}},F_D)$ itself. Its placement on the Karlsson Di\c{t}a circle at $\lambda\in\{\pi/2,\,3\pi/2\}$ (and nowhere else) rests on a \emph{numerical} CHM identification of $H(\theta_D,\pi/4,\lambda)$ with $D_0$/$D_{\mathrm{bc}}$ at those two $\lambda$ (dephased residual $\sim10^{-13}$; \repo{results/identify\_d0\_in\_karlsson.txt}). Neither point lies in $\Lambda_{\mathrm{cert}}$, so Theorems~\ref{thm:T3} and~\ref{thm:T4} cover disjoint point sets.

(ii) \emph{Clique scope.} Coverage is the single third basis $F_D$; this is \textbf{not} the all-pool-clique exhaustion of Theorem~\ref{thm:T3}, and it is not a theorem on the Di\c{t}a circle, on $\mathcal{R}$, or on all of $K_6^{(3)}$.

(iii) \emph{Object / method / result vs Brierley--Weigert (2009).} Headline: Theorem~\ref{thm:T4} \emph{reproves} a fact BW already established, by an independent exact method; it does \textbf{not} extend BW. Asymmetry: T4 is \emph{weaker in coverage} (one third basis $F_D$ vs BW's full pool $N_v{=}120$, $N_t{=}10$) and \emph{stronger in one narrow sense} (exact symbolic unit-ideal certificate for the fixed-$(H,B_3)$ witness $W_1$ over $\mathbb{Q}(\zeta_{24},\sqrt5)$, versus BW's Groebner-on-pool construction plus combinatorial $N_p{=}0$). $D_{\mathrm{bc}}$ is CHM-equivalent to textbook $D_0$, and $F_D$ is Bengtsson's closed-form third MUB---neither object is new. BW's $N_p{=}0$ already implies $W_1(D_0,B_3)=\emptyset$ for every third basis from the pool, so the mathematical result at this point is corollary-grade; the value here is a machine-checkable cross-check on that instance via a different ideal---a pipeline-validation result, not a co-equal headline with Theorems~\ref{thm:T1} and~\ref{thm:T3}. The same certificate at $\lambda=0$ or $\pi/3$ (points \emph{not} CHM-equivalent to $D_0$) remains open, blocked on an algebraic form of $B_3$ there.

(iv) \emph{Numerical counterpart (2026-09).} After the methodology hardening, the audited pipeline also reproduces BW's complete-pool counts at this very point ($N_v{=}120$, $N_t{=}10$, $N_p{=}0$, and $W_1$ certified empty for $10/10$ cliques; all counts certified-numerical with reseed crosschecks): T4's exact $\{1\}$ certificate and that numerical reproduction now cover the same instance by independent methods (exact Gr\"obner vs.\ certified homotopy), each cross-validating the other, with BW's original result remaining the prior for both.
\end{remark}


\begin{thebibliography}{99}

\bibitem{JaupiMethods2026}
A.~Jaupi, ``Computational Identification and Audit of Third-MUB Loci Within Karlsson's Three-Parameter Family,'' companion methods manuscript to the present theorems note, 2026.

\bibitem{Zauner1999}
G.~Zauner, ``Quantendesigns---Grundz\"uge einer nichtkommutativen Designtheorie,'' Ph.D. dissertation, Univ.\ Wien, 1999.

\bibitem{Scott2010}
A.~J.~Scott and M.~Grassl, ``SIC-POVMs and MUBs from algebraic number theory,'' \emph{J.\ Math.\ Phys.}, vol.~51, no.~4, p.~042203, 2010.

\bibitem{Bengtsson2007}
I.~Bengtsson, W.~Tadej, and K.~\.Zyczkowski, ``Five qubit entangled states, positive maps, and matrix pencils,'' \emph{Phys.\ Rev.\ A}, vol.~75, no.~3, p.~032319, 2007.

\bibitem{BengtssonMUH2007}
I.~Bengtsson, W.~Bruzda, {\AA}.~Ericsson, J.-{\AA}.~Larsson, W.~Tadej, and K.~\.Zyczkowski, ``Mutually unbiased bases and Hadamard matrices of order six,'' \emph{J.\ Math.\ Phys.}, vol.~48, no.~5, p.~052106, 2007. [Online]. Available: \url{https://arxiv.org/abs/quant-ph/0610161}

\bibitem{Brierley2009}
S.~Brierley and S.~Weigert, ``Constructing mutually unbiased bases in dimension six,'' \emph{Phys.\ Rev.\ A}, vol.~79, no.~5, p.~052316, 2009. [Online]. Available: \url{https://arxiv.org/abs/0901.4051}

\bibitem{Brierley2008}
S.~Brierley and S.~Weigert, ``Maximal sets of mutually unbiased quantum states in dimension six,'' \emph{Phys.\ Rev.\ A}, vol.~78, no.~4, p.~042312, 2008.

\bibitem{Grassl2004}
M.~Grassl, ``On SIC-POVMs and MUBs in dimension six,'' arXiv:quant-ph/0405174, 2004.

\bibitem{Ivanovic1981}
I.~D.~Ivanovi\'c, ``Geometrical description of quantal state determination,'' \emph{J.\ Phys.\ A}, vol.~14, no.~12, pp.~3241--3245, 1981.

\bibitem{Wootters1989}
W.~K.~Wootters and B.~D.~Fields, ``Optimal state-determination by mutually unbiased measurements,'' \emph{Ann.\ Phys.}, vol.~191, no.~2, pp.~363--381, 1989.

\bibitem{Karlsson2011}
A.~Karlsson, ``Three-parameter families of complex Hadamard matrices of order 6,'' arXiv:1003.4177, 2011.

\bibitem{Dita2009}
P.~Dita, ``Some results on the parametrization of complex Hadamard matrices of order 6,'' \emph{J.\ Phys.\ A}, vol.~42, no.~46, p.~465305, 2009.

\bibitem{McNulty2026}
D.~McNulty and S.~Weigert, ``Mutually Unbiased Bases in Composite Dimensions -- A Review,'' \emph{Quantum}, vol.~10, p.~2051, 2026.

\bibitem{HomotopyContinuation}
P.~Breiding and S.~Timme, ``HomotopyContinuation.jl: A package for homotopy continuation in Julia,'' in \emph{Proc.\ ICMS}, 2018. [Software] Version 2.22.

\end{thebibliography}

\end{document}
